% ============================================================
% Arquitectura Cliente-Servidor en Tres Capas — Cap. 4.1.1
% Proyecto BTH V2.0 — Sistema Web de Inscripción MTP San Julián
% Formato: TikZ vectorial editable (APA 7)
% Uso: % ============================================================
% Arquitectura Cliente-Servidor en Tres Capas — Cap. 4.1.1
% Proyecto BTH V2.0 — Sistema Web de Inscripción MTP San Julián
% Formato: TikZ vectorial editable (APA 7)
% Uso: % ============================================================
% Arquitectura Cliente-Servidor en Tres Capas — Cap. 4.1.1
% Proyecto BTH V2.0 — Sistema Web de Inscripción MTP San Julián
% Formato: TikZ vectorial editable (APA 7)
% Uso: % ============================================================
% Arquitectura Cliente-Servidor en Tres Capas — Cap. 4.1.1
% Proyecto BTH V2.0 — Sistema Web de Inscripción MTP San Julián
% Formato: TikZ vectorial editable (APA 7)
% Uso: \input{imagenes/arquitectura_tres_capas} dentro de figure
% ============================================================
\begin{tikzpicture}[
    capa/.style={
        rectangle, rounded corners=5pt,
        draw=primary!85, align=center,
        font=\small,
        minimum height=1.7cm,
        inner sep=6pt,
        text width=11.5cm
    },
    flecha/.style={
        -{Stealth[length=3mm,width=2.2mm]},
        thick, draw=primary!85
    },
    flechaInv/.style={
        {Stealth[length=3mm,width=2.2mm]}-,
        thick, draw=primary!60
    },
    etiquetaCapa/.style={
        font=\scriptsize\bfseries, text=primary, align=center
    },
    detalle/.style={
        font=\scriptsize, text=darktext, align=center
    },
    tituloCapa/.style={
        font=\small\bfseries, align=center
    }
]

% --- CAPA 1: PRESENTACIÓN (FRONTEND) ---
\node[capa, fill=lightgray, draw=primary!70] (frontend) at (0, 4.2) {
    \tituloCapa{1. Capa de Presentación \normalfont\textit{(Frontend)}}\\
    \detalle{HTML5 \textbullet{} CSS3 responsivo \textbullet{} JavaScript ES6}\\
    \detalle{\footnotesize Navegador del postulante (PC / móvil) — Validación reactiva en cliente}
};

\node[etiquetaCapa, rotate=90] at (-7.05, 4.2) {CLIENTE};

% Icono usuario sutil (texto)
\node[detalle, font=\tiny\itshape, text=primary!70] at (0, 5.32) {$\sim$400 postulantes/año \textbullet{} 7 especialidades \textbullet{} $<$1.5 MB/trámite};

% --- FLECHA FRONTEND ↔ BACKEND ---
\draw[flecha] (0, 3.30) -- (0, 2.30);
\node[detalle, font=\tiny, fill=white, inner sep=1pt] at (2.15, 2.80) {HTTPS / JSON};
\node[detalle, font=\tiny, fill=white, inner sep=1pt] at (-2.10, 2.80) {Validación};
\draw[flechaInv, dashed, draw=primary!45] (0, 2.30) -- (0, 3.30);

% --- CAPA 2: LÓGICA DE NEGOCIO (BACKEND) ---
\node[capa, fill=secondary!14, draw=primary!70] (backend) at (0, 1.2) {
    \tituloCapa{2. Capa de Lógica del Negocio \normalfont\textit{(Backend)}}\\
    \detalle{PHP 8.2 sobre Apache 2.4}\\
    \detalle{\footnotesize Sesiones, validación C.I.+fecha, reglas matrícula, agrupación paralelos}
};

\node[etiquetaCapa, rotate=90] at (-7.05, 1.2) {SERVIDOR};

% --- FLECHA BACKEND ↔ BD ---
\draw[flecha] (0, 0.30) -- (0, -0.70);
\node[detalle, font=\tiny, fill=white, inner sep=1pt] at (2.35, -0.20) {SQL transaccional};
\node[detalle, font=\tiny, fill=white, inner sep=1pt] at (-2.25, -0.20) {ACID};
\draw[flechaInv, dashed, draw=primary!45] (0, -0.70) -- (0, 0.30);

% --- CAPA 3: PERSISTENCIA (BASE DE DATOS) ---
\node[capa, fill=primary, text=white, draw=primary] (bd) at (0, -1.85) {
    {\small\bfseries 3. Capa de Persistencia \normalfont\textit{(Base de Datos)}}\\
    {\scriptsize PostgreSQL 16.4 — Modelo relacional normalizado}\\
    {\scriptsize\itshape Transacciones ACID + bloqueo pesimista \texttt{FOR UPDATE} — 0 sobrecupos}
};

\node[etiquetaCapa, rotate=90, text=white!90] at (-7.05, -1.85) {DATOS};

% --- Línea lateral vertical decorativa (opcional, sutil) ---
\draw[primary!25, line width=2.5pt, rounded corners=3pt] (-6.55, 5.55) -- (-6.55, -2.75);

% --- Nota de rendimiento ---
\node[detalle, font=\tiny\itshape, text=darktext!65, align=center] at (0, -2.95) {Validación OE4: $>$100 peticiones/min \textbullet{} latencia $<$1.2 s \textbullet{} usabilidad $\geq$90\% \textbullet{} errores $<$5\%};

% --- Borde sutil alrededor (desactivado para APA limpio) ---
% \begin{scope}[on background layer]
%   \node[fit=(frontend)(bd), draw=primary!12, rounded corners=7pt, inner sep=7pt] {};
% \end{scope}

\end{tikzpicture}
 dentro de figure
% ============================================================
\begin{tikzpicture}[
    capa/.style={
        rectangle, rounded corners=5pt,
        draw=primary!85, align=center,
        font=\small,
        minimum height=1.7cm,
        inner sep=6pt,
        text width=11.5cm
    },
    flecha/.style={
        -{Stealth[length=3mm,width=2.2mm]},
        thick, draw=primary!85
    },
    flechaInv/.style={
        {Stealth[length=3mm,width=2.2mm]}-,
        thick, draw=primary!60
    },
    etiquetaCapa/.style={
        font=\scriptsize\bfseries, text=primary, align=center
    },
    detalle/.style={
        font=\scriptsize, text=darktext, align=center
    },
    tituloCapa/.style={
        font=\small\bfseries, align=center
    }
]

% --- CAPA 1: PRESENTACIÓN (FRONTEND) ---
\node[capa, fill=lightgray, draw=primary!70] (frontend) at (0, 4.2) {
    \tituloCapa{1. Capa de Presentación \normalfont\textit{(Frontend)}}\\
    \detalle{HTML5 \textbullet{} CSS3 responsivo \textbullet{} JavaScript ES6}\\
    \detalle{\footnotesize Navegador del postulante (PC / móvil) — Validación reactiva en cliente}
};

\node[etiquetaCapa, rotate=90] at (-7.05, 4.2) {CLIENTE};

% Icono usuario sutil (texto)
\node[detalle, font=\tiny\itshape, text=primary!70] at (0, 5.32) {$\sim$400 postulantes/año \textbullet{} 7 especialidades \textbullet{} $<$1.5 MB/trámite};

% --- FLECHA FRONTEND ↔ BACKEND ---
\draw[flecha] (0, 3.30) -- (0, 2.30);
\node[detalle, font=\tiny, fill=white, inner sep=1pt] at (2.15, 2.80) {HTTPS / JSON};
\node[detalle, font=\tiny, fill=white, inner sep=1pt] at (-2.10, 2.80) {Validación};
\draw[flechaInv, dashed, draw=primary!45] (0, 2.30) -- (0, 3.30);

% --- CAPA 2: LÓGICA DE NEGOCIO (BACKEND) ---
\node[capa, fill=secondary!14, draw=primary!70] (backend) at (0, 1.2) {
    \tituloCapa{2. Capa de Lógica del Negocio \normalfont\textit{(Backend)}}\\
    \detalle{PHP 8.2 sobre Apache 2.4}\\
    \detalle{\footnotesize Sesiones, validación C.I.+fecha, reglas matrícula, agrupación paralelos}
};

\node[etiquetaCapa, rotate=90] at (-7.05, 1.2) {SERVIDOR};

% --- FLECHA BACKEND ↔ BD ---
\draw[flecha] (0, 0.30) -- (0, -0.70);
\node[detalle, font=\tiny, fill=white, inner sep=1pt] at (2.35, -0.20) {SQL transaccional};
\node[detalle, font=\tiny, fill=white, inner sep=1pt] at (-2.25, -0.20) {ACID};
\draw[flechaInv, dashed, draw=primary!45] (0, -0.70) -- (0, 0.30);

% --- CAPA 3: PERSISTENCIA (BASE DE DATOS) ---
\node[capa, fill=primary, text=white, draw=primary] (bd) at (0, -1.85) {
    {\small\bfseries 3. Capa de Persistencia \normalfont\textit{(Base de Datos)}}\\
    {\scriptsize PostgreSQL 16.4 — Modelo relacional normalizado}\\
    {\scriptsize\itshape Transacciones ACID + bloqueo pesimista \texttt{FOR UPDATE} — 0 sobrecupos}
};

\node[etiquetaCapa, rotate=90, text=white!90] at (-7.05, -1.85) {DATOS};

% --- Línea lateral vertical decorativa (opcional, sutil) ---
\draw[primary!25, line width=2.5pt, rounded corners=3pt] (-6.55, 5.55) -- (-6.55, -2.75);

% --- Nota de rendimiento ---
\node[detalle, font=\tiny\itshape, text=darktext!65, align=center] at (0, -2.95) {Validación OE4: $>$100 peticiones/min \textbullet{} latencia $<$1.2 s \textbullet{} usabilidad $\geq$90\% \textbullet{} errores $<$5\%};

% --- Borde sutil alrededor (desactivado para APA limpio) ---
% \begin{scope}[on background layer]
%   \node[fit=(frontend)(bd), draw=primary!12, rounded corners=7pt, inner sep=7pt] {};
% \end{scope}

\end{tikzpicture}
 dentro de figure
% ============================================================
\begin{tikzpicture}[
    capa/.style={
        rectangle, rounded corners=5pt,
        draw=primary!85, align=center,
        font=\small,
        minimum height=1.7cm,
        inner sep=6pt,
        text width=11.5cm
    },
    flecha/.style={
        -{Stealth[length=3mm,width=2.2mm]},
        thick, draw=primary!85
    },
    flechaInv/.style={
        {Stealth[length=3mm,width=2.2mm]}-,
        thick, draw=primary!60
    },
    etiquetaCapa/.style={
        font=\scriptsize\bfseries, text=primary, align=center
    },
    detalle/.style={
        font=\scriptsize, text=darktext, align=center
    },
    tituloCapa/.style={
        font=\small\bfseries, align=center
    }
]

% --- CAPA 1: PRESENTACIÓN (FRONTEND) ---
\node[capa, fill=lightgray, draw=primary!70] (frontend) at (0, 4.2) {
    \tituloCapa{1. Capa de Presentación \normalfont\textit{(Frontend)}}\\
    \detalle{HTML5 \textbullet{} CSS3 responsivo \textbullet{} JavaScript ES6}\\
    \detalle{\footnotesize Navegador del postulante (PC / móvil) — Validación reactiva en cliente}
};

\node[etiquetaCapa, rotate=90] at (-7.05, 4.2) {CLIENTE};

% Icono usuario sutil (texto)
\node[detalle, font=\tiny\itshape, text=primary!70] at (0, 5.32) {$\sim$400 postulantes/año \textbullet{} 7 especialidades \textbullet{} $<$1.5 MB/trámite};

% --- FLECHA FRONTEND ↔ BACKEND ---
\draw[flecha] (0, 3.30) -- (0, 2.30);
\node[detalle, font=\tiny, fill=white, inner sep=1pt] at (2.15, 2.80) {HTTPS / JSON};
\node[detalle, font=\tiny, fill=white, inner sep=1pt] at (-2.10, 2.80) {Validación};
\draw[flechaInv, dashed, draw=primary!45] (0, 2.30) -- (0, 3.30);

% --- CAPA 2: LÓGICA DE NEGOCIO (BACKEND) ---
\node[capa, fill=secondary!14, draw=primary!70] (backend) at (0, 1.2) {
    \tituloCapa{2. Capa de Lógica del Negocio \normalfont\textit{(Backend)}}\\
    \detalle{PHP 8.2 sobre Apache 2.4}\\
    \detalle{\footnotesize Sesiones, validación C.I.+fecha, reglas matrícula, agrupación paralelos}
};

\node[etiquetaCapa, rotate=90] at (-7.05, 1.2) {SERVIDOR};

% --- FLECHA BACKEND ↔ BD ---
\draw[flecha] (0, 0.30) -- (0, -0.70);
\node[detalle, font=\tiny, fill=white, inner sep=1pt] at (2.35, -0.20) {SQL transaccional};
\node[detalle, font=\tiny, fill=white, inner sep=1pt] at (-2.25, -0.20) {ACID};
\draw[flechaInv, dashed, draw=primary!45] (0, -0.70) -- (0, 0.30);

% --- CAPA 3: PERSISTENCIA (BASE DE DATOS) ---
\node[capa, fill=primary, text=white, draw=primary] (bd) at (0, -1.85) {
    {\small\bfseries 3. Capa de Persistencia \normalfont\textit{(Base de Datos)}}\\
    {\scriptsize PostgreSQL 16.4 — Modelo relacional normalizado}\\
    {\scriptsize\itshape Transacciones ACID + bloqueo pesimista \texttt{FOR UPDATE} — 0 sobrecupos}
};

\node[etiquetaCapa, rotate=90, text=white!90] at (-7.05, -1.85) {DATOS};

% --- Línea lateral vertical decorativa (opcional, sutil) ---
\draw[primary!25, line width=2.5pt, rounded corners=3pt] (-6.55, 5.55) -- (-6.55, -2.75);

% --- Nota de rendimiento ---
\node[detalle, font=\tiny\itshape, text=darktext!65, align=center] at (0, -2.95) {Validación OE4: $>$100 peticiones/min \textbullet{} latencia $<$1.2 s \textbullet{} usabilidad $\geq$90\% \textbullet{} errores $<$5\%};

% --- Borde sutil alrededor (desactivado para APA limpio) ---
% \begin{scope}[on background layer]
%   \node[fit=(frontend)(bd), draw=primary!12, rounded corners=7pt, inner sep=7pt] {};
% \end{scope}

\end{tikzpicture}
 dentro de figure
% ============================================================
\begin{tikzpicture}[
    capa/.style={
        rectangle, rounded corners=5pt,
        draw=primary!85, align=center,
        font=\small,
        minimum height=1.7cm,
        inner sep=6pt,
        text width=11.5cm
    },
    flecha/.style={
        -{Stealth[length=3mm,width=2.2mm]},
        thick, draw=primary!85
    },
    flechaInv/.style={
        {Stealth[length=3mm,width=2.2mm]}-,
        thick, draw=primary!60
    },
    etiquetaCapa/.style={
        font=\scriptsize\bfseries, text=primary, align=center
    },
    detalle/.style={
        font=\scriptsize, text=darktext, align=center
    },
    tituloCapa/.style={
        font=\small\bfseries, align=center
    }
]

% --- CAPA 1: PRESENTACIÓN (FRONTEND) ---
\node[capa, fill=lightgray, draw=primary!70] (frontend) at (0, 4.2) {
    \tituloCapa{1. Capa de Presentación \normalfont\textit{(Frontend)}}\\
    \detalle{HTML5 \textbullet{} CSS3 responsivo \textbullet{} JavaScript ES6}\\
    \detalle{\footnotesize Navegador del postulante (PC / móvil) — Validación reactiva en cliente}
};

\node[etiquetaCapa, rotate=90] at (-7.05, 4.2) {CLIENTE};

% Icono usuario sutil (texto)
\node[detalle, font=\tiny\itshape, text=primary!70] at (0, 5.32) {$\sim$400 postulantes/año \textbullet{} 7 especialidades \textbullet{} $<$1.5 MB/trámite};

% --- FLECHA FRONTEND ↔ BACKEND ---
\draw[flecha] (0, 3.30) -- (0, 2.30);
\node[detalle, font=\tiny, fill=white, inner sep=1pt] at (2.15, 2.80) {HTTPS / JSON};
\node[detalle, font=\tiny, fill=white, inner sep=1pt] at (-2.10, 2.80) {Validación};
\draw[flechaInv, dashed, draw=primary!45] (0, 2.30) -- (0, 3.30);

% --- CAPA 2: LÓGICA DE NEGOCIO (BACKEND) ---
\node[capa, fill=secondary!14, draw=primary!70] (backend) at (0, 1.2) {
    \tituloCapa{2. Capa de Lógica del Negocio \normalfont\textit{(Backend)}}\\
    \detalle{PHP 8.2 sobre Apache 2.4}\\
    \detalle{\footnotesize Sesiones, validación C.I.+fecha, reglas matrícula, agrupación paralelos}
};

\node[etiquetaCapa, rotate=90] at (-7.05, 1.2) {SERVIDOR};

% --- FLECHA BACKEND ↔ BD ---
\draw[flecha] (0, 0.30) -- (0, -0.70);
\node[detalle, font=\tiny, fill=white, inner sep=1pt] at (2.35, -0.20) {SQL transaccional};
\node[detalle, font=\tiny, fill=white, inner sep=1pt] at (-2.25, -0.20) {ACID};
\draw[flechaInv, dashed, draw=primary!45] (0, -0.70) -- (0, 0.30);

% --- CAPA 3: PERSISTENCIA (BASE DE DATOS) ---
\node[capa, fill=primary, text=white, draw=primary] (bd) at (0, -1.85) {
    {\small\bfseries 3. Capa de Persistencia \normalfont\textit{(Base de Datos)}}\\
    {\scriptsize PostgreSQL 16.4 — Modelo relacional normalizado}\\
    {\scriptsize\itshape Transacciones ACID + bloqueo pesimista \texttt{FOR UPDATE} — 0 sobrecupos}
};

\node[etiquetaCapa, rotate=90, text=white!90] at (-7.05, -1.85) {DATOS};

% --- Línea lateral vertical decorativa (opcional, sutil) ---
\draw[primary!25, line width=2.5pt, rounded corners=3pt] (-6.55, 5.55) -- (-6.55, -2.75);

% --- Nota de rendimiento ---
\node[detalle, font=\tiny\itshape, text=darktext!65, align=center] at (0, -2.95) {Validación OE4: $>$100 peticiones/min \textbullet{} latencia $<$1.2 s \textbullet{} usabilidad $\geq$90\% \textbullet{} errores $<$5\%};

% --- Borde sutil alrededor (desactivado para APA limpio) ---
% \begin{scope}[on background layer]
%   \node[fit=(frontend)(bd), draw=primary!12, rounded corners=7pt, inner sep=7pt] {};
% \end{scope}

\end{tikzpicture}
